\section{Theoretical Background}\label{sec:background}

\subsection{Large language models and their limits in medicine}

A large language model is a transformer network trained to predict the next token of text on a very large corpus and then further trained to follow instructions and hold a dialogue. Models such as GPT-4o \cite{openai2024gpt4o} and the open-weight Llama~3 family \cite{grattafiori2024llama3} can answer questions, summarise documents and produce structured output such as JSON. Singhal \emph{et al.} showed that instruction-tuned LLMs encode a substantial amount of clinical knowledge and can answer medical exam questions at a high level \cite{singhal2023}.

The same work, and a large literature on hallucination \cite{ji2023}, also shows the limits. An LLM generates the most plausible continuation of a text, not the most correct one. When its internal knowledge is incomplete it may produce a confident but false statement, and when asked for references it may produce citations that look real but do not exist \cite{walters2023}. The model also has no access to information outside its prompt, so it cannot know the data of a specific patient unless that data is given to it. For a clinical assistant these are serious problems: a fluent wrong answer is more dangerous than no answer, because it is more likely to be believed.

\subsection{Building domain experts from general LLMs}\label{sec:bg-expert}

There are three main ways to make a general LLM behave as an expert in a narrow field. They are complementary, and this thesis uses all three (Table~\ref{tab:approaches}).

\paragraph{Prompting.} The cheapest approach is a carefully written system prompt (``You are a clinical hematology expert\ldots''), together with rules and an output format. Prompting changes the style and focus of the answers, but it does not add knowledge the model lacks, it does not provide sources, and it does not prevent unsupported statements.

\paragraph{Retrieval-augmented generation (RAG).} RAG \cite{lewis2020} leaves the model unchanged and supplies relevant passages from a curated corpus at question time. The source documents are split into chunks, each chunk is converted into a vector by a sentence-embedding model \cite{reimers2019}, and the vectors are stored in a vector database. For each question, the chunks whose vectors are most similar to the question's vector are retrieved and placed in the prompt, together with an instruction to answer from them and to cite them. RAG has two important properties for a domain expert: the knowledge can be updated without retraining, and every answer can be checked against the passages it was based on. Comparative studies have found retrieval to be more reliable than fine-tuning for injecting new factual knowledge into an LLM \cite{ovadia2024}.

\paragraph{Domain fine-tuning.} Fine-tuning changes the model's weights using examples from the domain. Medical models such as Med-PaLM \cite{singhal2023} and Meditron \cite{chen2023meditron} were built this way, but full fine-tuning of a model with billions of parameters is too expensive for most projects. \emph{Parameter-efficient} methods solve this. LoRA \cite{hu2022lora} freezes the original weight matrices $W$ and learns a low-rank update, so that a layer computes
\begin{equation}
  h = W x + \frac{\alpha}{r}\, B A x, \qquad B \in \mathbb{R}^{d \times r},\ A \in \mathbb{R}^{r \times k},\ r \ll \min(d,k),
\end{equation}
where only $A$ and $B$ are trained. QLoRA \cite{dettmers2023qlora} additionally stores the frozen model in 4-bit precision, which makes it possible to fine-tune an 8-billion-parameter model on a single 16\,GB GPU such as the free NVIDIA T4 in Google Colab. Fine-tuning is particularly effective at teaching \emph{behaviour}: answer length and structure, terminology, citation format and cautious phrasing. It also produces an open model that can be run on local hardware, independent of a commercial API.

\paragraph{Synthetic training data.} When no suitable human-written training data exist, question--answer pairs can be generated by a stronger model from trusted source text, in the spirit of Self-Instruct \cite{wang2023selfinstruct}. This is fast and cheap, but the generated data inherit the generator's mistakes and style, which must be considered when the resulting model is evaluated.

\begin{table}[H]
\centering
\caption{Ways of turning a general LLM into a domain expert, and how this thesis uses them.}
\label{tab:approaches}
\small
\begin{tabularx}{\textwidth}{l X X X}
\toprule
 & \textbf{Adds knowledge?} & \textbf{Gives checkable sources?} & \textbf{Role in this thesis} \\
\midrule
Prompting & No & No & System prompts and JSON output schema \\
RAG & Yes, at query time & Yes, retrieved passages & Textbook knowledge base (Sec.~\ref{sec:kb}) \\
Fine-tuning (QLoRA) & Mostly style and behaviour & No (model may imitate a citation format) & Open, self-hostable expert (Sec.~\ref{sec:ft}) \\
Tools / agent & Yes, case data and reference values & Yes, logged tool calls & GPT-4o ReAct agent (Sec.~\ref{sec:agent}) \\
\bottomrule
\end{tabularx}
\end{table}

\subsection{Tool use and the ReAct agent paradigm}

Standard RAG performs one retrieval followed by one generation. ReAct \cite{yao2023react} generalises this by interleaving reasoning and acting: the model writes a thought, chooses an action (a call to a tool), observes the result, and repeats until it produces a final answer. Tools can be database searches, calculators or reference tables. For a domain expert this has two advantages. The model can look up exact reference values and case data instead of relying on memory, and every tool call is recorded, so the path to the final answer can be audited. LangGraph \cite{langgraph} provides an implementation of this loop with structured tool calling, which is used in this thesis.

\subsection{Evaluating a domain expert}

A domain expert can be evaluated in three complementary ways:
\begin{enumerate}
  \item \textbf{Multiple-choice questions} with known answers, such as the medical entrance-exam questions of MedMCQA \cite{pal2022medmcqa}. They give an objective accuracy, but test recognition of facts rather than explanation.
  \item \textbf{Similarity of free-text answers to reference answers}, measured by word overlap (ROUGE-L, based on the longest common subsequence \cite{lin2004rouge}) or by the cosine similarity of sentence embeddings. These measures are cheap and reproducible, but they reward answers that \emph{look like} the reference, not necessarily answers that are correct.
  \item \textbf{Graded judgements} of answers by clinicians or, more cheaply, by a strong LLM used as a judge \cite{zheng2023judge}. An LLM judge scales well but has its own biases and does not replace expert review.
\end{enumerate}
This thesis uses the first two, and inspects answers qualitatively (Section~\ref{sec:eval}).

\subsection{Automated peripheral blood-smear analysis}

Image-based analysers for blood smears have existed since the 2000s, first with classical image processing and later with deep learning. Acevedo \emph{et al.} showed that convolutional neural networks can classify single blood-cell images with very high accuracy \cite{acevedo2019} and published the PBC dataset of 17{,}092 cell images in eight classes, recorded with a CellaVision DM96 analyser \cite{acevedo2020}. For locating cells in a whole field of view, general object detectors of the YOLO family \cite{jocher2023yolov8} are fine-tuned on annotated smears such as BCCD \cite{mooney2018bccd} and the more recent TXL-PBC collection \cite{gan2024txlpbc}.

Most of this work stops at a label per cell; the interpretation of the differential as a whole is left to the human reader. Table~\ref{tab:related} places the present system among typical approaches. The comparison is qualitative: it describes capabilities of system families, not of every individual system.

\begin{table}[H]
\centering
\caption{Qualitative comparison of system families.}
\label{tab:related}
\footnotesize
\begin{tabularx}{\textwidth}{>{\raggedright\arraybackslash}p{2.6cm} X X X X}
\toprule
\textbf{System family} & \textbf{Image analysis} & \textbf{Uncertainty} & \textbf{Explanation} & \textbf{Clinical interpretation} \\
\midrule
General LLM chat (e.g.\ ChatGPT) & None (text only) & None & Free text, no sources & Generic, not case-specific \\
CNN cell classifiers & Pre-cropped cells & Usually softmax only & Rare or post-hoc & None \\
Detector + classifier pipelines & Detection and classification & Rarely passed on & Boxes, sometimes saliency & None \\
\textbf{This thesis} & YOLOv8s + EfficientNet-B0 & MC-Dropout level per cell, review flag & Boxes, Grad-CAM, citations, tool trace & RAG + tool-using LLM expert; QLoRA open model \\
\bottomrule
\end{tabularx}
\end{table}

\subsection{Uncertainty and explainability}

A neural network's softmax output is not a calibrated probability; modern networks tend to be over-confident \cite{guo2017}. Bayesian deep learning treats the network weights as random variables, but exact inference is intractable. Practical approximations include deep ensembles \cite{lakshminarayanan2017}, which train several networks and average them, and \emph{Monte Carlo Dropout} \cite{gal2016}, which keeps dropout active at prediction time and averages $T$ stochastic forward passes. MC-Dropout needs no retraining when the network already contains dropout layers, which is why it is used here; its disadvantage is that it is only an approximation and is usually less well calibrated than ensembles.

Grad-CAM \cite{selvaraju2017} produces a heatmap showing which image regions most influenced a class prediction. It weights the feature maps $A^k$ of a convolutional layer by the average gradient of the class score $y^c$:
\begin{equation}
  L^c_{\text{Grad-CAM}} = \mathrm{ReLU}\Big(\sum_k \alpha_k^c A^k\Big), \qquad \alpha_k^c = \frac{1}{Z}\sum_{i,j}\frac{\partial y^c}{\partial A^k_{ij}}.
\end{equation}
In a clinical setting Grad-CAM is most useful as a check: if the heatmap highlights the background or a neighbouring cell rather than the cell of interest, the prediction should be treated as suspect.

\subsection{Trustworthy AI in medical decision support}

Regulatory and policy documents --- the FDA's Good Machine Learning Practice principles \cite{fda2021gmlp}, the EU Artificial Intelligence Act \cite{euaiact2024} and the WHO guidance on large multi-modal models in health \cite{who2024lmm} --- converge on a few requirements for high-risk medical AI: transparency about inputs and limits, auditability of decisions, communication of uncertainty, and effective human oversight. These requirements guided the design of the system in this thesis. The system is a research prototype: it is not clinically validated, not a medical device, and not intended for use on patient data outside a research setting.

\clearpage
